\documentclass[a4paper,10pt]{article}
\usepackage[francais]{babel}
\usepackage{url}

\selectlanguage{francais}
../letter-ensiie.sty

\selectlanguage{francais}
\enCouleur

\begin{document}


\name{Dr. St\'ephane \textsc{Genaud}}
\fonction{Professeur en Informatique à l'ENSIIE\\Directeur du Campus de Strasbourg}
\telephoneillkirch{+33 (0)369 200 231}
\email{genaud@ensiie.fr}

\begin{letter}[Attestation du Directeur de Thèse - Etablissement d'Accueil]{à}{Qui de droit}

\vspace{-6mm}
  Je  soussigné, Professeur  (HDR)  Stéphane Genaud,  membre  du Laboratoire  de
  recherche ICube, à accueilli M. DIBAMOU MBEUYO Cyrille au laboratoire Icube du
  3 avril 2015 au 31 juillet 2015.  Je m’engage à l’accueillir à nouveau pour un
  séjour ponctuel  selon les besoins  de son  travail. J'atteste que  M. DIBAMOU
  MBEUYO Cyrille originaire  du Cameroun et inscrit à la  Faculté de Sciences de
  l'Université de Ngaoundéré, est engagé dans  des travaux de thèse sur le thème
  suivant : « Ordonnancement de la sauvegarde/reprise des tâches malléables dans
  un environnement virtuel de calcul haute performances ».

  M. DIBAMOU  MBEUYO Cyrille est  titulaire d’une Licence  et d’un Master  II en
  Informatique obtenu à la Faculté de Sciences de l'Université de Ngaoundéré. Ce
  qui lui  a permis  de s'inscrire en  thèse de Doctorat  afin de  poursuivre le
  travail de  recherche qu'il a commencé  en Mater. Il mène  ses recherches avec
  courage et détermination afin de parvenir aux résultats escomptés.

  La mobilité de M. DIBAMOU vers  l’université d'accueil lui permet de découvrir
  les avancées  et de nouer  des collaborations  dans son domaine  de recherche,
  celui  du  Cloud  Computing.  Ceci  lui permettra  de  mieux  appréhender  les
  différents concepts et techniques liés à sa thématique. Les résultats attendus
  permettront  de mettre  à  la  disposition des  chercheurs  de son  université
  d'origine,   une   plateforme   qui   leur   permettrait   de   conduire   des
  expérimentations.  Cette plateforme  qui pourra  être déployée  aux moyens  de
  techniques  de virtualisation,  viendra  en effet  répondre  aux problèmes  de
  manque d'infrastructures de type grille de calcul que rencontrent généralement
  les chercheurs des pays du Sud.  Dans le cadre de ses travaux de recherche, la
  collaboration  qu'il  entretient  avec  son université  d'accueil  lui  permet
  d'approfondir ses connaissances  dans son domaine et ainsi  de mieux concevoir
  et  analyser les  solutions  qui  répondent aux  problèmes  de son  Université
  d'origine.  Le parcours scolaire de M. DIBAMOU MBEUYO Cyrille a été sanctionné
  au secondaire, par l'obtention de son  diplôme de Baccalauréat série C dans un
  système francophone.  Et il  a poursuivi ses  études supérieures  en français,
  dans une université publique  de son pays. Il n'a de ce  fait de difficultés à
  mener ses travaux de recherche doctorale en français.

\end{letter}
\begin{flushright}
\includegraphics[width=.23\textwidth]{signgenaud.jpg}
\end{flushright}



\end{document}
